\begin{Definition}{HTML}
    -Definimos HTML como un lenguaje de estructura semántica de documento.
\end{Definition}

\begin{Definition}{DOM}
    -Document Object Model $\implies $ Desarrollo web \\
    Se define como el grafo $G = (V, E)$ de elementos visualizables en HTML
\end{Definition}

\textbf{Caso práctico}: Sea $(Y)$ y $S$ yo y el servidor respectivamente. Usando GET-messages 
envio la solicitud de que quiero mis mensajes. Luego el servidor debiese en su motor interno gestionar la 
solicitud y preparar el resultado ($S: JSON$). 
\\
\\ 
El DOM lo construye ($Y$) después de recibir $JSON(S)$.
\\ 
\\
Luego $Y: JSON \implies DOM \implies visualizar$
\\
\\
\underline{\textbf{Quien debe paginar}}: Definimos el paginar como el proceso usual de cortar el JSON. 
Debe hacerlo $(S)$. Podría hacerlo $(Y)$ pero no sería eficiente.

\begin{Definition}{STATEFUL}
    -Decimos que un servidor es STATEFUL si la tupla $(t_1, t_2, \dots, t_n)$
    no influye sobre la acción $A$ a ejecutar en $t_n$. Los tokens por ejemplo, no dependen de $t_{n-1}$.
    Buscar en la base de datos del server por diccionario/hash tampoco depende de $t_{n-k}$.
\end{Definition}

\begin{Definition}{REST}
    -Estilo arquitectonico donde manipulo recursos usando protocolos HTTP.\\
    Satisface GET, POST, PUT, DELETE. \\
    Es $Stateless$.
\end{Definition}

\begin{Definition}{Entidad}
    -Es un objeto del dominio. Tiene identidad propia y atributos.
\end{Definition}

\begin{Definition}{PUT}
    -Es un método de REST. Lo que hace es \textbf{reemplazar} un recurso existente.
\end{Definition}

\begin{Definition}{CURSOR}
    - Sea D un conjunto ordenado de datos. El cursor es un puntero a un dato especifico desde donde 
    continuamos la lectura. \\
    - Es un identificador de la posición en un conjunto de datos que permite
    continuar la lectura desde un punto (punto especifico, no indice). 
\end{Definition}

Sea $D$ tal que se modifica en LIFO-STACK. $D$ es ordenado. 
Definimos el cursor como $\{e_{j_\leq} \}$. Tal que si paso de $D$ a $D \cup e$
$$ \implies \{ e_{j_{\leq}} \}_{e \cup D} = \{ e_\leq \}_{D} $$

\begin{Definition}{OFFSET}
    -Lo definimos de igual forma que Cursor solo que es sobre un indice. Es decir, si rellenamos por LIFO y ejecutamos
    la accion va a devolver un set de datos distinto.
\end{Definition}

\begin{Lemma}{Pooling}
    -Supongamos que queremos mantener una conexion instantanea. 
    Un usuario necesita ver sus mensajes en vivo pese a que aun no llega ninguno. 
    Puede ejecutar el GET. El servidor lo recibe. Lo deja en WAIT. Cuando llege un mensaje el 
    servidor lo envia. (La conexion se puede dejar en standby)
\end{Lemma}

\begin{Definition}{WebSocket}
    -Protocolo  de comunicacion TCP persistente (en vivo por harto rato) que es bidireccional
\end{Definition}

\begin{Definition}{HEAP}
     - H es un arbol binario : H es min - heap $\iff$
    $$ \forall x \in H: x.sons \geq x $$
\end{Definition}
Por definicion observar el minimo es $\mathcal{O}(1)$. 
Extraer (es decir sacar y reordenar) es $\mathcal{O}(\log (n))$
Insertar a su vez es $\mathcal{O}(log (n))$. 
\begin{Lemma}{Insertion in HEAP}
    - Para insertar en (min-HEAP). Se inserta en la ultima posicion disponible. 
    Se usa el sig. algoritmo recursivo. 
        $$while: \; if \; not \; check(x.father, x) \then swap(x.father, x)$$ \\
    Es facil notar que la complejidad de esto es $\mathcal{\log(n)}$
\end{Lemma}

\subsection{Una aplicación de interés}
Supongamos que tenemos un servicio de mensajes tal que cada mensaje es $M = (id, timestamp)$. Un mensaje $j$ es mas reciente que uno $i$ si:
$$\iff t_i \leq t_j $$.
Nos gustaria tener una estructura para siempre tener los $k$ mensajes más recientes. Uno podría pensar en un STACK pero los mensajes no llegan en 
perfecto orden debido a la concurrencia. Lo que nos interesa es poder manejar eficientemente la inserción. Por lo tanto hacemos un MIN-HEAP(k). Es decir
en la raíz va a estar el más viejo. Luego al insertar el $t_j$ basta comparar con la raíz (si es que el HEAP esta lleno). Popear la raíz, hacer heapify e 
e insertar el $t_j$. Recordamos que las inserciones son $\mathcal{O}(\log(k))$. Si reordenasemos a cada rato el stack tendriamos $\mathcal{O}(k \log (k))$
\\
Ojo que construir un HEAP por heapify-downs es $\mathcal{O}(n)$

\subsection{Un problema de interés - Requests per minute}
Supongamos que tenemos mensajes (id, timestamp) que llegan de forma arbitraria. 
Queremos que el server maneje solo 100 requests por minuto. Como lo hacemos.
La solucion es una ventana continua. Es decir, en cada segundo el servidor contara hasta el minuto pasado 
y solo mantendra los que caen en esa ventana. El resto chao. Esa ventana va a ser de tamaño 100. 
\\
Para escalar basta tomar REDIS con operaciones atomicas. Es decir colapsar el conteo a una operacion atomica.

\begin{Definition}{LRU CACHE}
    - Definimos el LRU como Least Recently Used. Es decir un stack de K elementos tal que
    eliminamos la wea en mayor desuso. Tal que para cada atomo $a$ hay 2 metodos.
    GET(key) y PUT(key, value)
\end{Definition}
¿Como eliminamos lo que esta en mayor desuso? Con una lista doble ligada. 

\begin{Definition}{Lista doble ligada}
    - Es una colección de elementos tal que cada elemento tiene $x.before$ y $x.after$
\end{Definition}
Mover elementos entre si en la lista doble ligada es $O(1)$. Depende del elemento solamente.